
\section{Introduction}
\label{sec:intro}

An innocent little combinatorial counting problem
asks for the number of triangulations of a
finite grid of size $m\times n$.  That is, for $m,n\ge1$
we define $P_{m,n}:=\{0,1,\ldots,m\}\times\{0,1,\ldots,n\}$, ``the
grid''. Equivalently, the point configuration $P_{m,n}$ 
consists of all points of the integer lattice $\Z^2$
in the lattice rectangle $\conv(P_{m,n})=[0,m]\times[0,n]$
of area~$mn$.
Every triangulation of this rectangle point set that uses all the
points in $P_{m,n}$ has $(m+1)(n+1)=|P_{m,n}|$ vertices, 
$2mn$ facets/triangles,
and $3mn+m+n$ edges, $2(m+n)$ of them on the boundary, 
the other $3mn-m-n$ ones in the interior.
All the triangles are minimal lattice triangles of area $\frac12$
(that is, of determinant~$1$), which are referred to as
\emph{unimodular} triangles. The grid triangulations that
use all the points are thus called \emph{unimodular triangulations}.
The number of unimodular triangulations of the grid $P_{m,n}$
 will be denoted by $f(m,n)$.
 
 As an example, our first figure shows one unimodular triangulation of
 the $5\times 6$ grid:
\begin{center}
\input{a_triangulation3.pstex_t}
\end{center}
To get started, one notes that of course $f(m,n)=f(n,m)$,
one discovers with pleasure that $f(1,n)=\binom{2n}n$, and
one works out by hand with a bit of pain that $f(2,2)=64$ and $f(2,3)=852$.
Furthermore, one observes that the special triangulations that
decompose into $m$ vertical strips of width~$1$ yield the
lower bound
\begin{equation}
  \label{eq:easylow}
  f(m,n)\ \ge\ f(1,n)^m\ =\ \binom{2n}n^m\;
\end{equation}
so for larger $m$ and $n$ the numbers $f(m,n)$ get huge \emph{very soon};
for example, the bound yields $f(5,6) > 6 \cdot 10^{14}$.

It is equally interesting to study/enumerate more
general types of triangulations, such as
triangulations of finite point sets that do not necessarily use 
all the points, triangulations of general convex or non-convex
lattice polygons, triangulations of general position
point sets, triangulations of higher-dimensional point sets, etc.
None of these will appear here, but we refer interested readers to
Lee \cite{Lee} and De Loera, Rambau \& Santos \cite{DLRS}.

Lattice triangulations are basic combinatorial objects,
and they are fundamental discrete geometric structures;
so it is no surprise that they appear in various computational geometry
contexts. However, lattice triangulations have also been
studied intensively from different algebraic geometry angles. So, 
triangulations of a convex lattice polygon
\begin{compactitem}
\item[$\bullet$] provide
the data for Viro's \cite{Viro} famous construction method
of plane algebraic curves with prescribed combinatorics and topology,
related to Hilbert's sixteenth problem;
\item[$\bullet$]
appear in Gel'fand--Kapranov--Zelevinsky's \cite{GKZ} theory of
discriminants, where ``regular'' triangulations are in bijection with 
the vertices of the secondary polytope of the point configuration, and
\item[$\bullet$]
model torus-equivariant 
crepant resolution of singularities for toric three-folds,
where ``regular'' triangulations correspond to projective desingularizations;
see e.\,g., Kempf et~al.\ \cite[Chap.~3]{KKMS} and Dais \cite{Dais-3dToric}.
\end{compactitem}
The last two points pose the problem of counting or estimating the
number $\freg(m,n)$ of \emph{regular} triangulations of~$P_{m,n}$,
that is, of triangulations of $\conv(P_{m,n})$ whose triangles are the
domains of linearity for a piecewise linear convex function
(\emph{lifting function}); see also, for example, Sturmfels
\cite[Chap.~8]{Sturmfels-GroePoly}, Ziegler \cite[Lect.~5]{Zpoly}.
Motivated by the toric variety considerations, one would like to know
whether/that for large $m$ and~$n$ most triangulations are
non-regular.  There is no proof, but a lot of evidence that this
should be true (see, e.g., Table~\ref{tab:large}).

In this context, we must admit that despite the effort that has been
put into studying this question (see e.\,g.\ Hastings
\cite[Chap.~2]{Hastings}), it has not become clear what non-regular
triangulations really ``look like.'' The ``mother of all examples'' is
the whirlpool triangulation of the $3\times 3$ grid:
\begin{center}
  \begin{overpic}[scale=2]{whirlpool.eps}
    \put(37,29){$p_0$}
    \put(54,38){$p_1$}
    \put(54,68){$p_2$}
    \put(37,60){$p_3$}
    \put(-8,3){$q_0$}
    \put(100,3){$q_1$}
    \put(100,93){$q_2$}
    \put(-8,93){$q_3$}
  \end{overpic}
\end{center}

Indeed, this unimodular triangulation is irregular, since a lifting
function~$h$ would have to satisfy
$h(p_i)+h(q_{i-1})<h(p_{i-1})+h(q_i)$ for $i=0,\dots,3$ (all indices
taken modulo~$4$), implying the contradiction $h(q_0)+\dots+
h(q_3)<h(q_0)+\dots+ h(q_3)$.

For $m=n=3$, except for three symmetric copies
this is the only example of a non-regular triangulation: We
have $f(3,3)-\freg(3,3)=4$.
This may lead one to the conjecture that some kind of
``generalized whirlpools'' are responsible for non-regularity.
However, the pictures of ``non-regular triangulations''
that we present below do not support this intuition.
\medskip

The plan for this paper is as follows:
In Section~\ref{sec:values} we face the challenge to
\emph{\bfseries count} explicitly, trying to cope with the
``combinatorial explosion.''
For this, we present a simple dynamic programming technique by which we
get surprisingly far, and which on the specific problem of
grid triangulations outruns the much more
sophisticated general techniques such as 
Avis \& Fukuda's \cite{AvisFukuda} reverse search, 
Aichholzer's \cite{Aichholzer-path} path of a triangulation,
or the oriented matroid technique of Rambau \cite{Rambau-TOPCOM}.

In Section~\ref{sec:bounds} we \emph{\bfseries estimate} $f(m,n)$ for
large $m$ and $n$, trying to narrow the bounds for the asymptotics.
It is interesting to compare with the situation for $N=(m+1)(n+1)$
points in the plane in general position, where the currently best
available upper bound seems to be Santos \& Seidel's
\cite{SantosSeidel} estimate that there are $o(59^N)$ triangulations.
However, in our problem the $N$ points are not at all in general
position --- and the upper bounds that we have to offer are much
better: We report on a neat $O(2^{3N})$ upper bound by Anclin
\cite{Anclin}, which substantially improves on a previous $O(2^{4N})$
upper bound by Orevkov \cite{Orevkov}.  Based on explicit enumeration
results from Section~\ref{sec:values}, we get a lower bound of
$2^{2.055\,mn}$ when both $m$ and~$n$ get large; note
that~(\ref{eq:easylow}) yields already a lower bound
$2^{(1-\littleo(1))2mn}$.

Finally, in Section~\ref{sec:trias} 
we \emph{\bfseries sample} lattice triangulations for large parameters
$m$ and~$n$,
and thus try to understand what typical lattice triangulations,
as well as typical regular lattice triangulations, ``look like.''
While we have some pictures to offer, proofs seem harder to come by.
Indeed, the pictures display some long-range order;
while this may make lattice triangulations interesting as a
statistical physics model, it generates serious obstacles for any proof that
the obvious Markov chain is rapidly mixing, and thus to
application of the (by now) standard theory \cite{Behrends-MarkovChains}.


%%% Local Variables:
%%% mode: latex
%%% TeX-master: "bcc_ct"
%%% End:
